\section{Discussion}\label{sec:discussion}

\subsection{Symmetry decides what can be reconstructed}

The reconstruction tasks were designed from the symmetry of the model rather than from the capacity of the learners, and the outcome is correspondingly sharp. Delay coordinates of the collar height cannot determine the sign of the spring force, because the reflection $\theta\to-\theta$ leaves the collar height invariant, and the original FPRM, which regresses on delay coordinates alone, therefore returns the average of the two branches, with a correlation near zero ($|\rho|<0.05$) at every missing fraction. This is an instance of a general principle of observability theory: the quality of a reconstruction depends on the observable, and symmetries can render an observable blind to part of the state~\cite{LetellierAguirre2002,Letellier2005}. The FPRM inherits this constraint from Takens' theorem, and the cross-mapping test of Fig.~\ref{fig:fprm_scheme} offers a cheap way to detect it before any model is trained. Two pieces of information restore identifiability in the PIAR: the observed neighbors of the incomplete channel, through continuity, and the phase of the forcing, which distinguishes mirror states of the driven system. A practical rule follows for the design of sensor layouts: a channel that is invariant under a symmetry of the dynamics should not be the only complete record for a channel that the symmetry changes, unless another signal, such as a drive or a reference clock, resolves that ambiguity.

\subsection{Distributions rather than points}

When the ambiguity is genuine, a conditional generative model is the appropriate tool, because it represents the branches instead of averaging them. Point metrics computed on the posterior median can then misrank reconstructions, as the higher correlation but worse CRPS of the cVAE in task C2 illustrates, and proper scoring rules~\cite{Gneiting2007}, together with the fraction of samples on the correct branch, are the relevant measures. Distributional fidelity also matters downstream: the invariant measure of the attractor, probed here through the Poincar\'e section, is recovered by the flow and blurred by the original FPRM (Fig.~\ref{fig:reconstruction}). The choice of engine can therefore be guided by the cross-mapping test. When the test reveals a mirror branch, a conditional generative model is required, unless observed neighbors resolve the branch, as in task C2 at small missing fractions. When the test indicates identifiability, a Gaussian process on an informative input is a strong and inexpensive baseline, provided that its training set resolves the folded attractor, as in task C3 with the phase but not in task C1; the nearest-neighbor library test of Sec.~\ref{sec:res_c1} checks this condition at negligible cost. Among generative models, a flow is preferable when physical residuals are to be imposed or many samples are required, and a diffusion model when labels are scarce and no residuals are imposed, or when conservative intervals matter more than computational cost (Sec.~\ref{sec:res_engine}).

\subsection{Physics as a double-edged sword}

Equation-of-motion residuals helped exactly where they carried information about the unknown degree of freedom: in task C1, the gravitational and forcing terms discriminate the sign of $\cos\theta$, and the kinematic relation fixes the angular velocity. Where the residuals were nearly invariant under the ambiguity, or dominated by the noise of numerical derivatives, they did not add information about the branch and, at $60$--$90\%$ missing data, pushed the flow to collapse onto a single branch. Three practical rules follow: a residual should be imposed only when it is informative about the ambiguous degree of freedom and well conditioned with respect to the noise of its derivatives; regularized models should be validated on imputations of artificially masked observed entries rather than on the likelihood of labeled ones; and sharpened posteriors must be checked for calibration (Sec.~\ref{sec:res_engine}). The topological origin of the collapse suggests remedies that preserve multimodality by construction: mixture base distributions, continuously indexed flows~\cite{Cornish2020}, and flows defined on the circle~\cite{Rezende2020}, which respect the periodicity of the angle. Structure-preserving models offer a complementary route to physical consistency. Hamiltonian and Lagrangian networks~\cite{Greydanus2019,Cranmer2020} and their dissipative extensions~\cite{Sosanya2022} could supply the dynamical prior of a future FPRM through the architecture rather than through the loss; the Hamiltonians learned by generic networks are not separable, and explicit symplectic integrators for nonseparable Hamiltonians~\cite{Tao2016} would then be needed to roll them out.

\subsection{Limitations and outlook}

The study is confined to a single degree of freedom, to a single chaotic reference state, to a harmonic torque rather than to the multiplicative base excitation of real isolators, and to synthetic data with missing-completely-at-random masks and a fixed noise level; gaps that are correlated in time, or caused by saturation, would require the missingness mechanism to be modeled~\cite{LittleRubin2019}. The benchmark uses two random masks per setting, which suffices to separate the qualitative regimes reported here but not to rank imputers whose scores differ by a few percent. The original FPRM was reimplemented with an exact Gaussian process; scalable approximations would allow it to use every labeled row at small missing fractions.

Several extensions follow naturally. An experimental PIAR instrumented with a load cell and a linear encoder would test the framework on real sensor records, and chains of coupled pendulums, in which one site is fully observed, would test it in higher dimension. Reconstructed records could also feed system identification, from the sparse regression of governing equations~\cite{Brunton2016} to the Bayesian estimation of physical parameters. For the latter, spectral Whittle likelihoods~\cite{Whittle1953,Sykulski2019} can be sampled with ensemble Markov chain Monte Carlo~\cite{GoodmanWeare2010,ForemanMackey2013} or with nested sampling~\cite{Skilling2006,Speagle2020}, which also yields the evidence for comparing competing models through Bayes factors~\cite{KassRaftery1995}, and the calibration of the resulting posteriors can be checked by simulation-based calibration~\cite{Talts2018}. Propagating the imputation uncertainty of the flow into such inferences by multiple imputation~\cite{LittleRubin2019} is a direct benefit of a generative engine that a point regressor cannot provide.

Finally, the same laboratory can be used in teaching. The five conservative regimes illustrate the classification of orbits by energy, the degenerate critical point, and the lengthening of the period predicted by the normal form; the driven system displays the Feigenbaum cascade, the Melnikov mechanism, and a strange attractor; and the FPRM turns Takens' theorem into a concrete, testable computation. Because every figure and number can be regenerated, the PIAR can accompany a course from analytical mechanics to data-driven dynamics.
